\section{Step 4: Transport the basis without losing the Wronskian}
\label{app:connection-matrices}
% BEGIN CURATED SUBJECT INDEX
\index{spectrum!point}
\index{Wronskian}
\index{connection coefficient}
\index{Stokes curve}
\index{monodromy}
% END CURATED SUBJECT INDEX

Let $\bm\Psi=(\psi_+,\psi_-)^{\mathsf T}$ be a local WKB basis with constant
Wronskian.  Crossing a simple Stokes curve multiplies its coefficient vector by
one of the triangular matrices in Eq.~\eqref{eq:stokes-matrices}.  Transport
between reference points $a$ and $b$ gives a diagonal matrix
\begin{equation}
  P_{ab}
  =\begin{pmatrix}
  \rme^{\VorosPeriod_{ab}}&0\\
  0&\rme^{-\VorosPeriod_{ab}}
  \end{pmatrix},
  \qquad
  \VorosPeriod_{ab}=\int_a^b\vsub{S}{odd}\dd x .
  \label{eq:propagation-matrix}
\end{equation}
A global connection matrix is an ordered product
\begin{equation}
  M_{\Gamma}
  =M_{j_n}P_{n-1,n}\cdots M_{j_2}P_{12}M_{j_1} .
  \label{eq:global-connection-product}
\end{equation}
Changing local normalization conjugates or diagonally rescales
$M_{\Gamma}$.  Its relevant zero condition and monodromy eigenvalues are
unchanged.

For an equation without a first-derivative term, the Wronskian
\begin{equation}
  W[\psi_1,\psi_2]
  =\psi_1\partial_x\psi_2
  -\psi_2\partial_x\psi_1
  \label{eq:wronskian}
\end{equation}
is independent of $x$.  Consequently every exact connection matrix has unit
determinant after consistent basis normalization.  This algebraic constraint
is a useful check on symbolic and numerical Stokes calculations.
